\SourcePage{165}
\SGAsection{2}{阿贝尔层的双扩张概念}
\textbf{2.0。} 本节以下始终以 $P,Q,G$ 表示\Term{topos}{拓扑斯} $\mathcal T$ 的\Term{commutative-group}{交换群}。我们将经常使用\Term{cartesian-diagram}{笛卡儿图}
\[
\begin{tikzcd}
P\arrow[d]&P\times Q\arrow[l]\arrow[d]\\
e&Q\arrow[l]
\end{tikzcd}
\tag{2.0.1}
\]
以及 $P$（相应地，$Q$、$P\times Q$）上的\Term{group}{群} $G_P$（相应地，$G_Q$、$G_{P\times Q}$），它们是 $e$ 上的\Term{group}{群} $G$ 的\Term{inverse-image}{逆像}。我们将 $G_{P\times Q}$ 与 $G_P$ 沿 $\operatorname{pr}_1$ 的\Term{inverse-image}{逆像} $(G_P)_{P\times Q}$ 或与 $G_Q$ 沿 $\operatorname{pr}_2$ 的\Term{inverse-image}{逆像} $(G_Q)_{P\times Q}$ 等同，二者均可。在这些记号中，$P,Q,P\times Q$ 的\Term{group}{群}结构不起\Term{action}{作用}。事实上，我们几乎从不需要使用 $P\times Q$ 上的积\Term{group}{群}结构。相反，考虑 $P\times Q$ 上由 $e$ 上的 $Q$ 的结构取\Term{inverse-image}{逆像}所得的 $P$-\Term{group}{群}结构，将是有用的：
\[
P\times Q=Q_P；\tag{2.0.2}
\]
同样还有由 $e$ 上的 $P$ 的结构取\Term{inverse-image}{逆像}所得的 $Q$-\Term{group}{群}结构：
\[
P\times Q=P_Q。\tag{2.0.3}
\]
这样，在 $P$ 上要考虑两个\Term{group}{群} $G_P$ 和 $Q_P$，在 $Q$ 上则要考虑两个\Term{group}{群} $G_Q$ 和 $P_Q$。若 $E$ 是 $P\times Q$ 上结构\Term{group}{群}为 $G_{P\times Q}$ 的一个\Term{torsor}{挠子}，则既可把它看作 $Q_P$ 上结构\Term{group}{群}为 $G_{P\times Q}=(G_P)_{Q_P}$ 的\Term{torsor}{挠子}，也可把它看作 $P_Q$ 上结构\Term{group}{群}为 $G_{P\times Q}=(G_Q)_{P_Q}$ 的\Term{torsor}{挠子}。因此，在 $E$ 上考虑一个与其\Term{torsor}{挠子}结构相容的、$Q_P$ 以 $G_P$ 为核的（交换）\Term{group-extension}{扩张}结构，或\Term{symmetric}{对称}地，
\SourcePage{166}
考虑一个与其\Term{torsor}{挠子}结构相容的、$P_Q$ 以 $G_Q$ 为核的\Term{group-extension}{扩张}结构，都是有意义的。此外，把 1.6 应用于\Term{induced-topos}{诱导拓扑斯} $\mathcal T/P$ 和 $\mathcal T/Q$，便得到具体描述这些结构的一种方便方法。利用\Term{canonical-isomorphism}{典范同构}
\[
\mathcal T/(P\times Q)\simeq(\mathcal T/P)/(P\times Q),\qquad
\mathcal T/(P\times Q)\simeq(\mathcal T/Q)/(P\times Q),
\]
可得如下描述。在\Term{torsor}{挠子} $E$ 上给定一个 $Q_P$ 以 $G_P$ 为核的\Term{group-extension}{扩张}结构，\Term{equivalence}{等价}于对 $\mathcal T$ 的任意\Term{object}{对象} $S$ 和任意元素 $p,p'\in P(S)$、$q\in Q(S)$，给定一个 $G_S$-\Term{torsor}{挠子}\Term{isomorphism}{同构}
\[
\varphi_{p,p';q}:E_{p,q}E_{p',q}\xrightarrow{\sim}E_{pp',q},\tag{2.0.4}
\]
并满足下面重述的条件。像通常一样（参见 1.1.2.1），对于任意点 $(p,q)\in(P\times Q)(S)=P(S)\times Q(S)$，以 $E_{p,q}$ 表示 $G_{P\times Q}$-\Term{torsor}{挠子} $E$ 沿 $(p,q):S\to P\times Q$ 的\Term{inverse-image}{逆像}，它是一个 $G_S$-\Term{torsor}{挠子}。须对数据 (2.0.4) 施加的条件是：关于 $S$ 的\Term{functoriality}{函子性}（即与 $\mathcal T$ 中\Term{morphism}{态射} $S'\to S$ 所给出的\Term{inverse-image}{逆像}相容）；\Term{associativity}{结合性}，它由形如 (1.1.4.1) 的图的\Term{commutativity}{交换性}表达：
\[
\begin{tikzcd}[column sep=small]
&E_{pp',q}E_{p'',q}\arrow[dr,"{\varphi_{pp',p'';q}}"]&\\
E_{p,q}E_{p',q}E_{p'',q}\arrow[ur,"{\varphi_{p,p';q}\wedge\mathrm{id}}"]\arrow[dr,"{\mathrm{id}\wedge\varphi_{p',p'';q}}"']&&E_{pp'p'',q}\\
&E_{p,q}E_{p'p'',q}\arrow[ur,"{\varphi_{p,p'p'';q}}"']&
\end{tikzcd}
\tag{2.0.5}
\]
其中 $p,p',p''\in P(S)$、$q\in Q(S)$；最后是\Term{commutativity}{交换性}，它由形如 (1.2.1) 的图的\Term{commutativity}{交换性}表达：
\SourcePage{167}
\[
\begin{tikzcd}
E_{p,q}E_{p',q}\arrow[r,"{\varphi_{p,p';q}}"]\arrow[d,"\mathrm{sym}"']&E_{pp',q}\arrow[d,"\mathrm{id}"]\\
E_{p',q}E_{p,q}\arrow[r,"{\varphi_{p',p;q}}"']&E_{p'p,q}
\end{tikzcd}
\tag{2.0.6}
\]
其中 $p,p'\in P(S)$、$q\in Q(S)$。同样，可用 $G_S$-\Term{torsor}{挠子}\Term{isomorphism}{同构}
\[
\psi_{p;q,q'}:E_{p,q}E_{p,q'}\xrightarrow{\sim}E_{pp',q},\tag{2.0.7}
\]
具体描述 $E$ 上一个 $P_Q$ 以 $G_Q$ 为核的\Term{group-extension}{扩张}结构；这里 $p\in P(S)$、$q,q'\in Q(S)$，$\psi$ 关于 $S$ 具有\Term{functoriality}{函子性}，并满足与 (2.0.5)、(2.0.6) \Term{symmetric}{对称}的\Term{associativity}{结合性}与\Term{commutativity}{交换性}条件，即下图交换：
\[
\begin{tikzcd}[column sep=small]
&E_{p,qq'}E_{p,q''}\arrow[dr,"{\psi_{p;qq',q''}}"]&\\
E_{p,q}E_{p,q'}E_{p,q''}\arrow[ur,"{\psi_{p;q,q'}\wedge\mathrm{id}}"]\arrow[dr,"{\mathrm{id}\wedge\psi_{p;q',q''}}"']&&E_{p,qq'q''}\\
&E_{p,q}E_{p,q'q''}\arrow[ur,"{\psi_{p;q,q'q''}}"']&
\end{tikzcd}
\tag{2.0.8}
\]
其中 $p\in P(S)$、$q,q',q''\in Q(S)$，以及
\[
\begin{tikzcd}
E_{p,q}E_{p,q'}\arrow[r,"{\psi_{p;q,q'}}"]\arrow[d,"\mathrm{sym}"']&E_{p,qq'}\arrow[d,"\mathrm{id}"]\\
E_{p,q'}E_{p,q}\arrow[r,"{\psi_{p;q',q}}"']&E_{p,q'q}
\end{tikzcd}
\tag{2.0.9}
\]
其中 $p\in P(S)$、$q,q'\in Q(S)$。

\SourcePage{168}
\begin{SGAresult}{定义}{2.1}
设 $P,Q,G$ 是 $\mathcal T$ 的\Term{abelian-group}{阿贝尔群}，$E$ 是一个 $G_{P\times Q}$-\Term{torsor}{挠子}，并赋有由\Term{isomorphism}{同构} (2.0.4) 定义的、$Q_P$ 以 $G_P$ 为核的\Term{group-extension}{扩张}结构（这些\Term{isomorphism}{同构}关于 $S$ 具有\Term{functoriality}{函子性}，且服从由图 (2.0.5)、(2.0.6) 的\Term{commutativity}{交换性}表达的条件），以及由\Term{isomorphism}{同构} (2.0.7) 定义的、$P_Q$ 以 $G_Q$ 为核的\Term{group-extension}{扩张}结构（这些\Term{isomorphism}{同构}关于 $S$ 具有\Term{functoriality}{函子性}，且服从由图 (2.0.8)、(2.0.9) 的\Term{commutativity}{交换性}表达的条件）。若对于 $\mathcal T$ 的任意\Term{object}{对象} $S$ 和任意 $p,p'\in P(S)$、$q,q'\in Q(S)$，下图交换，就称上述两个结构相容：
\[
\begin{tikzcd}[column sep=small,row sep=large]
&E_{p,qq'}E_{p',qq'}\arrow[ddr,"{\varphi_{p,p';qq'}}"]&\\
E_{p,q}E_{p,q'}E_{p',q}E_{p',q'}\arrow[ur,"{\psi_{p;q,q'}\wedge\psi_{p';q,q'}}"]\arrow[d]&&\\
E_{p,q}E_{p',q}E_{p,q'}E_{p',q'}\arrow[dr,"{\varphi_{p,p';q}\wedge\varphi_{p,p';q'}}"']&&E_{pp',qq'}\\
&E_{pp',q}E_{pp',q'}\arrow[ur,"{\psi_{pp';q,q'}}"']&
\end{tikzcd}
\tag{2.1.1}
\]
其中竖直\Term{arrow}{箭头}由\Term{symmetric}{对称}\Term{isomorphism}{同构}
\[
E_{p,q'}E_{p',q}\xrightarrow{\sim}E_{p',q}E_{p,q'}
\]
\Term{derive}{导出}。所谓 $(P,Q)$ 以 $G$ 为核的\Term{biextension}{双扩张}，是指一个 $G_{P\times Q}$-\Term{torsor}{挠子} $E$，赋有一个 $Q_P$ 以 $G_P$ 为核的\Term{group-extension}{扩张}结构，以及一个 $G_Q$ 以 $P_Q$ 为核的\Term{group-extension}{扩张}结构，并且这两个结构在上述意义下相容。
\end{SGAresult}

\SourcePage{169}
\textbf{2.1.1。} 也可以说，在 $G_{P\times Q}$-\Term{torsor}{挠子} $E$ 上，一个\Term{biextension}{双扩张}结构由 $E$ 上两个并非处处有定义的\Term{composition-law}{合成律} $\varphi(x,y)$ 和 $\psi(x,y)$ 给出；第一个在 $\operatorname{pr}_2j(x)=\operatorname{pr}_2j(y)$ 时有定义，第二个在 $\operatorname{pr}_1j(x)=\operatorname{pr}_1j(y)$ 时有定义（其中 $j:E\to P\times Q$ 是\Term{torsor}{挠子}的结构\Term{morphism}{态射}）。除了形如
\[
\varphi(gx,y)=\varphi(x,gy)=g\varphi(x,y)\qquad(x,y\in E(S),\ g\in G(S))\tag{2.1.1.1}
\]
的\Term{associativity}{结合性}及 $\psi$ 的类似关系外，公理便由这些部分\Term{composition-law}{合成律}的\Term{associativity}{结合性}和\Term{commutativity}{交换性}，以及关系
\[
\varphi\bigl(\psi(x,y),\psi(z,t)\bigr)=\psi\bigl(\varphi(x,z),\varphi(y,t)\bigr)\tag{2.1.1.2}
\]
来表达；这里 $x,y,z,t$ 分别属于 $E_{p,q}(S),E_{p,q'}(S),E_{p',q}(S),E_{p',q'}(S)$。

\textbf{2.2。\Term{unit-section}{单位截面}。} 我们应考虑把 $E$ 看作 $Q_P$ 以 $G_P$ 为核的\Term{group-extension}{扩张}时的\Term{unit-section}{单位截面}
\[
e_{E/P}\in\Gamma(E/P)=\operatorname{Hom}_P(P,E),\tag{2.2.1}
\]
以及\Term{symmetric}{对称}地，把 $E$ 看作 $P_G$ 以 $G_Q$ 为核的\Term{group-extension}{扩张}时的\Term{unit-section}{单位截面}
\[
e_{E/Q}\in\Gamma(E/Q)=\operatorname{Hom}_Q(Q,E)。\tag{2.2.2}
\]
特别地，$P\times Q$ 上的\Term{torsor}{挠子} $E$（结构\Term{group}{群}为 $G_{P\times Q}$），在 $1_P\times Q$ 上由 $e_{E/P}$ \Term{canonical}{典范}地\Term{trivialization}{平凡化}，在 $P\times1_Q$ 上由 $e_{E/Q}$ \Term{canonical}{典范}地\Term{trivialization}{平凡化}。此外，这些\Term{trivialization}{平凡化}在 $P\times Q$ 的\Term{unit-section}{单位截面} $1_P\times1_Q$ 上重合，因而定义同一个\Term{section}{截面}
\SourcePage{170}
\[
e_E\in\Gamma(E),\tag{2.2.3}
\]
我们稍稍滥用语言，称之为 $E$ 的\Term{unit-section}{单位截面}（或中性\Term{section}{截面}）。为证明所说的、分别由 $e_{E/P}$ 和 $e_{E/Q}$ \Term{derive}{导出}的 $E$ 的两个\Term{section}{截面} $e,e'$ 相等，注意它们分别刻画为 $P\times Q$ 的\Term{unit-section}{单位截面}在 $E$ 中的\Term{inverse-image}{逆像} $E_{1_P,1_Q}$ 关于两个\Term{composition-law}{合成律} $\varphi$ 和 $\psi$ 的\Term{unit-section}{单位截面}。现将 (2.1.1.2) 应用于 $E_{1,1}$ 的\Term{section}{截面} $e,e',e',e$，便发现两边分别等于 $e$ 和 $e'$，从而如所断言的那样有 $e=e'$。还可进一步\Term{pushout}{推出}，$E_{1,1}$ 上的两个\Term{composition-law}{合成律} $\varphi$ 和 $\psi$ 相同。的确，显然它们分别借助\Term{morphism}{态射} $g\mapsto g.e$ 和 $g\mapsto g.e'$ 传递结构，都对应于 $G$ 的\Term{group}{群}律；因 $e=e'$，故它们相同。

还应指出，在底 $P$ 上 $Q_P$ 以 $G_P$ 为核的\Term{group-extension}{扩张} $E$，\Term{restriction}{限制}到 $P$ 的\Term{section}{截面} $1_P$ 上，便在 $1_P\times Q$ 于 $G$ 中的\Term{inverse-image}{逆像}上给出一个 $Q$ 以 $G$ 为核的\Term{group-extension}{扩张}结构。立即可见，\Term{homomorphism}{同态} $e_{E/Q}$ 对这个\Term{group-extension}{扩张}结构是\Term{multiplication}{乘法}\Term{homomorphism}{同态}，故这个\Term{group-extension}{扩张}是\Term{canonical}{典范}\Term{split-extension}{分裂}的。\Term{symmetric}{对称}地，把 $Q$ 上 $P_Q$ 以 $G_Q$ 为核的\Term{group-extension}{扩张} $E$ \Term{restriction}{限制}到 $Q$ 的\Term{section}{截面} $1_Q$ 上所得到的 $P$ 以 $G$ 为核的\Term{group-extension}{扩张}，由 $e_{E/P}$ 给出\Term{canonical}{典范}\Term{split-extension}{分裂}。

\SourcePage{171}
\textbf{2.3。\Term{biextension}{双扩张}的\Term{morphism}{态射}。} 设 $E$ 是 $(P,Q)$ 以 $G$ 为核的一个\Term{biextension}{双扩张}，同样，$E'$ 是 $(P',Q')$ 以 $G'$ 为核的一个\Term{biextension}{双扩张}，其中 $P',Q',G'$ 仍是 $\mathcal T$ 的\Term{commutative-group}{交换群}。所谓从 $E$ 到 $E'$ 的\Term{biextension}{双扩张}\Term{morphism}{态射}，是指一组数据 $(u,v,w,f)$，其中
\[
u:P\longrightarrow P',\qquad v:Q\longrightarrow Q',\qquad w:G\longrightarrow G'\tag{2.3.1}
\]
是\Term{group}{群}\Term{morphism}{态射}，而
\[
f:E\longrightarrow E'\tag{2.3.2}
\]
是 $E,E'$ 的底层\Term{set}{集合}\Term{sheaf}{层}之间的\Term{morphism}{态射}，使得 $f$ 既是与相对\Term{group}{群}\Term{homomorphism}{同态}
\[
u\times v:Q_P\longrightarrow Q'_{P'},\qquad u\times w:G_P\longrightarrow G'_{P'}
\]
关联的、分别在 $P$ 和 $P'$ 上的\Term{group-extension}{群扩张}\Term{homomorphism}{同态}，又是与相对\Term{group}{群}\Term{homomorphism}{同态}
\[
u\times v:P_Q\longrightarrow P'_{Q'},\qquad v\times w:G_Q\longrightarrow G'_{Q'}
\]
关联的、分别在 $Q$ 和 $Q'$ 上的\Term{group-extension}{群扩张}\Term{homomorphism}{同态}。还应指出，仅知道 $f$ 就能知道 $u,v,w$；因为从 $f$ 取商可恢复 $u\times v$，再分别\Term{restriction}{限制}到 $P\times1_Q$ 和 $1_P\times Q$，便恢复 $u$ 和 $v$，而把 $f$ \Term{restriction}{限制}到 $G.e_E$ 就恢复 $w$。所以，把\Term{morphism}{态射} $(u,v,w,f)$ 与 $f$ 等同起来并无不妥。

只要像通常那样，小心地把源和靶附加到\Term{morphism}{态射}上（行一个\Term{canonical}{典范}的屈膝礼），便可看到，\Term{topos}{拓扑斯} $\mathcal T$ 的\Term{biextension}{双扩张}构成一个\Term{category}{范畴}，其\Term{object}{对象}和\Term{morphism}{态射}由 2.1 和 2.3 定义，而\Term{morphism}{态射}的复合
\SourcePage{172}
按显然的方式定义。此外，通过 $(P,Q,G,E)\mapsto(P,Q,G)$，得到一个\Term{canonical}{典范}\Term{functor}{函子}
\[
\operatorname{BIEXT}(\mathcal T)\longrightarrow\operatorname{GRAB}(\mathcal T)\times\operatorname{GRAB}(\mathcal T)\times\operatorname{GRAB}(\mathcal T),\tag{2.3.3}
\]
其中 $\operatorname{BIEXT}(\mathcal T)$ 表示 $\mathcal T$ 上的\Term{biextension}{双扩张}\Term{category}{范畴}，$\operatorname{GRAB}(\mathcal T)$ 表示 $\mathcal T$ 上的\Term{abelian-group}{阿贝尔群}\Term{category}{范畴}。还可将上述\Term{canonical}{典范}\Term{functor}{函子}分成两个，即 $E\mapsto(P,Q)$：
\[
\operatorname{BIEXT}(\mathcal T)\longrightarrow\operatorname{GRAB}(\mathcal T)\times\operatorname{GRAB}(\mathcal T)\tag{2.3.4}
\]
和 $E\mapsto G$：
\[
\operatorname{BIEXT}(\mathcal T)\longrightarrow\operatorname{GRAB}(\mathcal T)。\tag{2.3.5}
\]

\textbf{2.4。$P,Q,G$ 可变时的\Term{category}{范畴} $\operatorname{BIEXT}(P,Q,G)$。} 我断言，\Term{canonical}{典范}\Term{functor}{函子} (2.3.4) 是\Term{fibered-functor}{纤维化的}，而\Term{canonical}{典范}\Term{functor}{函子} (2.3.5) 是\Term{cofibered-functor}{余纤维化的}。因而第一个断言特别意味着，对 $(P,Q)$ 以 $G$ 为核的一个\Term{biextension}{双扩张} $E$，以及\Term{abelian-group}{阿贝尔群}\Term{morphism}{态射}
\[
u:P'\longrightarrow P,\qquad v:Q'\longrightarrow Q,
\]
可找到 $(P',Q')$ 以 $G$ 为核的一个\Term{biextension}{双扩张} $E'$，以及一个与 $u,v,\mathrm{id}_G$ 相容的\Term{biextension}{双扩张}\Term{morphism}{态射}
\[
f:E'\longrightarrow E,
\]
它对于进入 $E$ 且与 $u,v,\mathrm{id}_G$ 相容的\Term{biextension}{双扩张}\Term{morphism}{态射}，具有通常的\Term{universal-property}{泛性质}。第二个断言也可同样展开，特别意味着：对 $(P,G)$ 以 $G$ 为核的一个\Term{biextension}{双扩张} $E$，以及一个\Term{abelian-group}{阿贝尔群}\Term{morphism}{态射}
\[
w:G\longrightarrow G',
\]
\SourcePage{173}
存在 $(P,Q)$ 以 $G'$ 为核的一个\Term{biextension}{双扩张} $E'$，以及一个与 $\mathrm{id}_P,\mathrm{id}_Q,w$ 相容的\Term{biextension}{双扩张}\Term{morphism}{态射}
\[
g:E\longrightarrow E',
\]
它对于从 $E$ 出发且与 $\mathrm{id}_P,\mathrm{id}_Q,w$ 相容的\Term{biextension}{双扩张}\Term{morphism}{态射}，具有通常的\Term{universal-property}{泛性质}。这里断言存在的两个构造，留给读者具体写出；只须照搬通常的\Term{group-extension}{群扩张}情形中熟知的类似构造，便很容易完成。

因此，对 $\mathcal T$ 的三个\Term{abelian-group}{阿贝尔群} $P,Q,G$，以
\[
\operatorname{BIEXT}(P,Q;G)\tag{2.4.1}
\]
表示 $(P,Q)$ 以 $G$ 为核的\Term{biextension}{双扩张}\Term{category}{范畴}，定义为 (2.3.3) 在 $(P,Q,G)$ 处的\Term{fiber-category}{纤维范畴}（即其\Term{morphism}{态射}取为那些与 $P,Q,G$ 的恒等\Term{morphism}{态射}相容的\Term{biextension}{双扩张}\Term{morphism}{态射}），便可看到，这个\Term{category}{范畴}“\Term{covariant}{协变地}依赖于 $G$，\Term{contravariant}{反变地}依赖于 $P,Q$”。此外，按照上述\Term{variance}{变性}，\Term{category}{范畴} $\operatorname{BIEXT}(P,Q;G)$ 对三个变元 $P,Q,G$ 中的每一个都具有\Term{additivity}{加性}。例如，对于变元 $G$，这一断言的确切含义如下：

 a) $\operatorname{BIEXT}(P,Q;0)$ 与单点\Term{category}{范畴}\Term{equivalence-of-categories}{等价}；

 b) 对 $\mathcal T$ 的两个\Term{abelian-group}{阿贝尔群} $G,G'$，由两个\Term{projection}{投影} $G\times G'\to G$ 和 $G\times G'\to G'$ \Term{derive}{导出}的\Term{functor}{函子}
\[
\operatorname{BIEXT}(P,G;G\times G')\longrightarrow
\operatorname{BIEXT}(P,Q;G)\times\operatorname{BIEXT}(P,Q;G')\tag{2.4.2}
\]
是一个\Term{equivalence-of-categories}{范畴等价}。用 [4] 的语言也可以说，对固定的 $P,Q$，$(P,Q)$ 以可变\Term{group}{群} $G$ 为核的\Term{biextension}{双扩张}\Term{category}{范畴}，
\SourcePage{174}
经由 (2.3.5) 诱导的\Term{functor}{函子} $E\mapsto G$，是 $\operatorname{GRAB}(\mathcal T)$ 上的一个\Term{additivity}{加性}\Term{cofibered-category}{余纤维范畴}。

\textbf{注记 2.4.3。} 按上述文献的术语，前述\Term{cofibered-category}{余纤维范畴}不仅具有\Term{additivity}{加性}，甚至还是“\Term{left-exact}{左正合}”的 [4, 2.2]。这里不验证这个性质；下文（3.6）将从更精确的事实\Term{pushout}{推出}它。这里只指出，上述文献的\Term{generic}{一般}理论将使我们能够用 $\operatorname{GRAB}(\mathcal T)$ 中的一个\Term{chain-complex}{链复形} $L_\bullet$（只依赖于 $P$ 和 $Q$），在\Term{fibered-category}{纤维范畴}\Term{equivalence-of-categories}{等价}的意义下具体描述这个\Term{cofibered-category}{余纤维范畴}的结构。下文也将给出这个\Term{complex}{复形} $L_\bullet$ 的一个直接构造，即\Term{complex}{复形} $P\mathbin{\overset{\mathrm L}{\otimes}}Q$（\Term{tensor-product}{张量积}取在一个\Term{derived-category}{导出范畴}中）。

\textbf{2.5。\Term{category}{范畴} $\operatorname{BIEXT}(P,Q;G)$ 的结构。} [4, 第 1 节] 的\Term{generic}{一般}\Term{form}{形式}考察，使我们可以从\Term{category}{范畴} $\operatorname{BIEXT}(P,Q;G)$ \Term{additivity}{加性地}依赖于 $G$ 这一事实，\Term{pushout}{推出}一些有用的结论。

\textbf{2.5.1。} 首先，固定 $G$ 时，\Term{category}{范畴} $\operatorname{BIEXT}(P,Q;G)$ 赋有一个合成结构，它以具有\Term{functoriality}{函子性}的方式，把 $(P,Q)$ 以 $G$ 为核的两个\Term{biextension}{双扩张} $E,E'$ 对应到一个乘积\Term{biextension}{双扩张}，记作 $E.E'$。可以认为，这个构造不过是照搬一个\Term{commutative-group}{交换群}以另一个\Term{commutative-group}{交换群}为核的两个\Term{group-extension}{扩张}的乘积这一熟知构造；还应指出，把 $E$ 和 $E'$ 解释为 $Q_P$ 以 $G_P$ 为核的\Term{group-extension}{群扩张}，或解释为 $P_Q$ 以 $G_Q$ 为核的\Term{group-extension}{群扩张}时，这个构造与后一概念确实相容。特别地，$EE'$ 作为 $G_{P\times Q}$-\Term{torsor}{挠子}，
\SourcePage{175}
就是由 $E,E'$ 所定义的两个 $G_{P\times Q}$-\Term{torsor}{挠子}的乘积。此外，$E.E'$ 的形成在相差\Term{canonical-isomorphism}{典范同构}的意义下具有\Term{associativity}{结合性}和\Term{commutativity}{交换性}；$\operatorname{BIEXT}(P,Q;G)$ 中所考虑的\Term{composition-law}{合成律}有一个“双侧中性\Term{object}{对象}”，称为 $(P,Q)$ 以 $G$ 为核的\Term{trivial-biextension}{平凡双扩张}。它\Term{constructible}{可构造}为（同上）由 $(P,Q)$ 以零\Term{group}{群} $0$ 为核的唯一\Term{biextension}{双扩张} $E_0$（相差\Term{isomorphism}{同构}），借助\Term{homomorphism}{同态} $0\to G$ 所\Term{derive}{导出}的\Term{biextension}{双扩张}。根据这个描述立即可验证，\Term{trivial-biextension}{平凡双扩张} $E$ 也可描述为\Term{trivial-torsor}{平凡挠子}
\[
E=P\times Q\times G,
\]
其在 $P$ 上和 $Q$ 上的两个\Term{group}{群}结构，分别由 $Q\times G$ 和 $P\times G$ 的积\Term{group}{群}结构，沿 $P\to e$ 和 $Q\to e$ 作\Term{base-change}{基变换}而得。

最后，对 $(P,Q)$ 以 $G$ 为核的任意\Term{biextension}{双扩张} $E$，可定义其\Term{inverse-biextension}{逆双扩张}，记作 $E^{-1}$，并赋有 $EE^{-1}$ 与\Term{trivial-biextension}{平凡双扩张}之间的一个\Term{isomorphism}{同构}；这在相差唯一\Term{isomorphism}{同构}的意义下确定 $E^{-1}$，并使我们能够定义它对于 $E$ 的\Term{functoriality}{函子性}依赖。

总之，\Term{category}{范畴} $\operatorname{BIEXT}(P,Q;G)$ 看来类似于\Term{category}{范畴}的\Term{category}{范畴} $(\mathrm{Cat})$ 中的一个\Term{group-object}{群对象}，只是\Term{associativity}{结合性}、\Term{commutativity}{交换性}以及\Term{unit}{单位}和逆元关系，仅在给定的\Term{isomorphism}{同构}之下成立。我们可以说，这个\Term{category}{范畴}是（$2$-\Term{category}{范畴} $(\mathrm{Cat})$ 的）一个\Term{pseudo-group}{伪群}。

\SourcePage{176}
\textbf{2.5.2。} 固定 $P,Q,G$，由上述性质可得以下事实。对 $(P,Q)$ 以 $G$ 为核的任意\Term{biextension}{双扩张} $E$，其\Term{endomorphism}{自态射}\Term{monoid}{幺半群}与\Term{trivial-biextension}{平凡双扩张} $\mathbf1$ 的\Term{endomorphism}{自态射}\Term{monoid}{幺半群}\Term{canonical}{典范}\Term{isomorphic}{同构}：给后者的每个\Term{endomorphism}{自态射} $u$ 对应 $\mathbf1.E\simeq E$ 的\Term{endomorphism}{自态射} $u.\mathrm{id}_E$ 即可。另一方面，$\mathbf1$ 的\Term{endomorphism}{自态射}\Term{monoid}{幺半群}是一个\Term{commutative-group}{交换群}，记作
\[
\operatorname{Biext}^0(P,Q;G)；\tag{2.5.2.1}
\]
其\Term{composition-law}{合成律}也由
\[
(u,v)\longmapsto u.v\in\operatorname{End}(\mathbf1.\mathbf1)\simeq\operatorname{End}(\mathbf1)
\]
给出。最后，\Term{composition-law}{合成律} $E,E'\mapsto EE'$ 在 $(P,Q)$ 以 $G$ 为核的\Term{biextension}{双扩张}的\Term{isomorphism}{同构}类\Term{set}{集合}
\[
\operatorname{Biext}^1(P,Q;G)\tag{2.5.2.2}
\]
中定义一个\Term{commutative-group}{交换群}律；$E$ 的类的逆元是\Term{inverse-bitorsor}{逆双挠子} $E^{-1}$ 的类。

\textbf{2.5.3。} 给定一个\Term{group}{群}\Term{homomorphism}{同态} $G\to G'$ 时，相应的\Term{functor}{函子}
\[
\operatorname{BIEXT}(P,Q;G)\longrightarrow\operatorname{BIEXT}(P,Q;G')\tag{2.5.3.1}
\]
在相差\Term{canonical-isomorphism}{典范同构}的意义下，与 2.5.1 中考虑的运算 $E.E'$ 和 $E^{-1}$ 相容。特别地，由此得出，两个显然的映射
\[
\operatorname{Biext}^i(P,Q;G)\longrightarrow\operatorname{Biext}^i(P,Q;G')\qquad(i=0,1)\tag{2.5.3.2}
\]
是\Term{group}{群}\Term{homomorphism}{同态}。此外，若 $E$ 是 $(P,Q)$ 以 $G$ 为核的一个\Term{biextension}{双扩张}，$E'$ 是相应的 $(P,Q)$ 以 $G'$ 为核的\Term{biextension}{双扩张}，则自然\Term{homomorphism}{同态}
\[
\operatorname{Hom}(E,E)=\operatorname{Aut}(E)\longrightarrow
\operatorname{Hom}(E',E')=\operatorname{Aut}(E')\tag{2.5.3.3}
\]
正是 $i=0$ 时的 (2.5.3.2)。

\SourcePage{177}
\textbf{2.5.4。} 还可容易地具体写出\Term{group}{群} $\operatorname{Biext}^0(P,Q;G)$。的确，若 $E$ 是 $(P,Q)$ 以 $G$ 为核的一个\Term{biextension}{双扩张}，则 $E$ 的\Term{automorphism}{自同构} $u$ 首先是 $E$ 的底层 $G_{P\times Q}$-\Term{torsor}{挠子}的\Term{automorphism}{自同构}，因此由\Term{morphism}{态射}
\[
u_0:P\times Q\longrightarrow G
\]
通过公式
\[
u(x)=u_0(j(x))x
\]
给出。在此基础上，立即可见，$u$ 与 $E$ 上第一个（相应地，第二个）\Term{group-extension}{扩张}结构相容，当且仅当 $u_0$ 对第二个（相应地，第一个）变元是\Term{additivity}{加性的}。因此，$u$ 是 $E$ 的\Term{biextension}{双扩张}结构的一个\Term{automorphism}{自同构}，当且仅当 $u_0$ 是\Term{bilinear}{双线性}的。于是得到一个\Term{canonical}{典范}\Term{bijection}{双射}
\[
\operatorname{Biext}^0(P,Q;G)\xrightarrow{\sim}\operatorname{Hom}(P\otimes Q,G)。\tag{2.5.4.1}
\]
容易验证，它与 $E$ 无关，保持\Term{group}{群}结构，并且关于 $G$ 具有\Term{functoriality}{函子性}（第一项的\Term{functoriality}{函子性}由 (2.5.3.2) 具体给出）。

下文将确定\Term{group}{群} $\operatorname{Biext}^1(P,Q;G)$，并将发现它\Term{canonical}{典范}地、且关于 $G$ 具有\Term{functoriality}{函子性}地\Term{isomorphic}{同构}于
\[
\operatorname{Ext}^1\bigl(P\mathbin{\overset{\mathrm L}{\otimes}}Q,G\bigr)。
\]

\textbf{2.6。} 在 2.5 中，我们只让 $G$ 变化，而保持 $P,Q$ 固定。现在应简要谈谈 $P,Q$ 的变化。其余事实都将由之\Term{form}{形式}地\Term{pushout}{推出}的基本事实是，$\operatorname{BIEXT}(P,Q;G)$ 关于 $P,Q$ 的\Term{contravariant}{反变性}，与它关于 $G$ 的\Term{covariant}{协变性}可交换，因此所考虑的\Term{category}{范畴}可看作一个\Term{cofibered-category}{余纤维范畴}的\Term{fiber-category}{纤维范畴}：
\SourcePage{178}
\[
\operatorname{BIEXT}'(\mathcal T)\longrightarrow
\operatorname{GRAB}(\mathcal T)^\circ\times\operatorname{GRAB}(\mathcal T)^\circ\times\operatorname{GRAB}(\mathcal T),\tag{2.6.1}
\]
其中上标 $\circ$ 表示\Term{opposite-category}{反范畴}（所考虑的结构\Term{functor}{函子}为 $E\mapsto(P,Q,G)$）。由此，所考虑的\Term{cofibered-category}{余纤维范畴}对这三个变元中的每一个都有\Term{additivity}{加性}，或者按我们的说法，它是\Term{triadditive}{三重加性的}。照搬一些标准论证——这里不再重述——便得出，\Term{fiber-category}{纤维范畴} $\operatorname{BIEXT}(P,P';G)$ 上的合成结构 $E.E'$，可以用三个变元中任一个的\Term{additivity}{加性}结构来描述，效果相同（即所得的三个\Term{pseudo-group}{伪群}结构\Term{canonical}{典范}\Term{isomorphic}{同构}）。由此又可得出，\Term{abelian-group}{阿贝尔群}
\[
\operatorname{Biext}^i(P,Q;G),\qquad i=0,1,\tag{2.6.2}
\]
可看作关于取值于 $\operatorname{GRAB}(\mathcal T)^\circ$、$\operatorname{GRAB}(\mathcal T)^\circ$ 和 $\operatorname{GRAB}(\mathcal T)$ 的变元 $P,Q,G$ 的两个\Term{trilinear}{三线性}\Term{functor}{函子}；或者，也可看作关于取值于 $\operatorname{GRAB}(\mathcal T)$ 的可变 $P,Q,G$ 的\Term{trilinear}{三线性}\Term{functor}{函子}，关于 $P,Q$ \Term{contravariant}{反变}，关于 $G$ \Term{covariant}{协变}，取值于\Term{abelian-group}{阿贝尔群}\Term{category}{范畴} $\operatorname{GRAB}$。还可验证，对于这样明确的\Term{variance}{变性}，\Term{isomorphism}{同构} (2.5.4.1) 关于 $P,Q,G$ 都具有\Term{functoriality}{函子性}，而不仅是关于 $G$。下文用一个 $\operatorname{Ext}^1$ 来确定 $\operatorname{Biext}^1$ 时，也将如此。

\textbf{2.7。给定\Term{biextension}{双扩张}的\Term{symmetric-biextension}{对称双扩张}。} 若 $E$ 是 $(P,Q)$ 以 $G$ 为核的一个\Term{biextension}{双扩张}，则可定义 $E$ 的\Term{symmetric-biextension}{对称双扩张} ${}^sE$，它是 $(Q,P)$ 以 $G$ 为核的一个\Term{biextension}{双扩张}。${}^sE$ 的底层\Term{set}{集合}\Term{sheaf}{层}就是 $E$ 的底层\Term{set}{集合}\Term{sheaf}{层}；$G$ 在 $E$ 和 ${}^sE$ 上的\Term{group-action}{作用}相同；结构\Term{morphism}{态射} ${}^sE\to P\times Q$ 是结构\Term{morphism}{态射} $E\to P\times Q$ 与\Term{symmetric}{对称} $P\times Q\to Q\times P$ 的复合；最后，
\SourcePage{179}
${}^sE$ 的左右\Term{group-extension}{扩张}结构，分别是 $E$ 的右左\Term{group-extension}{扩张}结构。显然，运算 $E\mapsto{}^sE$ 是\Term{biextension}{双扩张}\Term{category}{范畴} $\operatorname{BIEXT}(\mathcal T)$ 的一个\Term{automorphism}{自同构}，而且事实上是一个\Term{involution}{对合}\Term{automorphism}{自同构}：
\[
{}^s({}^sE)=E。
\]
对于关于\Term{bitorsor}{双挠子} $E,\ldots$ 的任何命题，把它应用于 ${}^sE,\ldots$，便可\Term{pushout}{推出}一个\Term{symmetric}{对称}的命题。我们可称之为“交换左右”而得的命题。当然，\Term{generic}{一般}只须具体写出两个\Term{symmetric}{对称}命题之一；至于选左还是选右，不妨按个人的政治偏好来定。

\textbf{2.8。\Term{topos}{拓扑斯}变换。\Term{localization}{局部化}。} 设给定一个\Term{topos}{拓扑斯}\Term{morphism}{态射}
\[
f:\mathcal T'\longrightarrow\mathcal T。
\]
那么，“\Term{inverse-image-functor}{逆像函子}” $f^*$ 把 $(P,Q)$ 以 $G$ 为核的一个\Term{biextension}{双扩张}变成 $(f^*(P),f^*(Q))$ 以 $f^*(G)$ 为核的一个\Term{biextension}{双扩张} $f^*(E)$，正如每当某种结构可“用\Term{finite}{有限}\Term{projective-limit}{投射极限}和任意\Term{inductive-limit}{归纳极限}来表述”时一样。为证明 $f^*(E)$ 上由 $E$ 的两个\Term{group-extension}{扩张}结构\Term{derive}{导出}的两个\Term{group-extension}{扩张}结构相容，可用“泛”情形中的图 (2.1.1) 的\Term{commutativity}{交换性}表达这一相容性，即取 $S=P\times P\times Q\times Q$，且 $p,p',q,q'$ 为 $S$ 到这四个因子的四个\Term{projection}{投影}。照例，通过 $E\mapsto f^*(E)$，得到一个\Term{canonical}{典范}\Term{functor}{函子}
\[
f^*:\operatorname{BIEXT}(\mathcal T)\longrightarrow\operatorname{BIEXT}(\mathcal T')。
\]
\SourcePage{180}
以下我们将直接使用而不再说明：借助\Term{biextension}{双扩张}所作的所有构造，都与\Term{topos}{拓扑斯}\Term{morphism}{态射}的\Term{inverse-image}{逆像}相容。特别地，2.5 中考虑的乘积运算 $EE'$ 和\Term{inversion}{取逆}运算 $E^{-1}$ 都是如此。

特别可取 $\mathcal T'$ 为一个\Term{induced-topos}{诱导拓扑斯} $\mathcal T/S'$，取 $f$ 为\Term{canonical}{典范}\Term{morphism}{态射}（称为“\Term{localization}{局部化}”\Term{morphism}{态射}，SGA 4 第 IV 讲 5）
\[
f=i_{S'}:\mathcal T/S'\longrightarrow\mathcal T。
\]
这时称 $f^*(E)$ 为 $E$ 在 $S'$ 上诱导的\Term{biextension}{双扩张}，也像通常那样记作 $E|S'$ 或 $E_{S'}$。由于\Term{restriction}{限制}运算显然具有通常的\Term{transitivity}{传递性}质，因此当 $S$ 为 $\mathcal T$ 的可变\Term{object}{对象}时，各\Term{category}{范畴} $\operatorname{BIEXT}(\mathcal T/S)$ 构成 $\mathcal T$ 上一个\Term{fibered-category}{纤维范畴}的\Term{fiber}{纤维}（SGA 1 第 VI 讲 8）。事实上，容易验证这个\Term{category}{范畴}甚至是 Giraud 意义下 $\mathcal T$ 上的一个\Term{stack}{叠} [2]，即“这个\Term{fibered-category}{纤维范畴}的\Term{fiber}{纤维}中的\Term{morphism}{态射}和\Term{object}{对象}都可以\Term{gluing}{粘合}”。例如，要在相差\Term{canonical-isomorphism}{典范同构}的意义下给定 $S$ 上一个\Term{biextension}{双扩张}，只须给定一个\Term{covering-family}{覆盖族}（即\Term{epimorphism}{满态射}族）$S_i\to S$，对每个 $i\in I$ 给定 $S_i$ 上（即 $\mathcal T/S_i$ 中）的一个\Term{biextension}{双扩张} $E_i$，对每对 $(i,j)$ 给定 $E_{iS_j}$ 与 $E_{jS_i}$ 的一个\Term{isomorphism}{同构}，并令这些\Term{isomorphism}{同构}满足通常的\Term{cocycle-condition}{余圈条件}。我们把这样定义的\Term{stack}{叠}称为 $\mathcal T$ 上的\Term{biextension}{双扩张}\Term{stack}{叠}。当用 $\mathcal T/S$ 代替 $\mathcal T$ 并令 $S$ 变化时，\Term{functor}{函子} (2.3.3)，从而\Term{derive}{导出}的\Term{functor}{函子} (2.3.4)、(2.3.5)，对应于 $\mathcal T$ 上的\Term{stack}{叠}\Term{morphism}{态射}。对可变的 $S$，\Term{category}{范畴} $\operatorname{BIEXT}(P_S,Q_S;G_S)$ 中的运算 $E.E'$ 和 $E^{-1}$（2.5.1）也是如此，这里 $P,Q,G$ 是 $\mathcal T$ 上给定的\Term{commutative-group}{交换群}。

\SourcePage{181}
\textbf{2.9。与\Term{birigidification}{双刚化}\Term{torsor}{挠子}的关系。}

采用 1.3.2 的术语，我们已在 2.2 中看到，对 $(P,Q)$ 以 $G$ 为核的任意\Term{biextension}{双扩张} $E$，其底层 $G_{P\times Q}$-\Term{torsor}{挠子}结构，相对于由 $P$ 和 $Q$ 的\Term{unit-section}{单位截面}组成的那一对，具有\Term{canonical}{典范}的\Term{birigidification}{双刚化}。于是得到一个\Term{canonical}{典范}\Term{functor}{函子}
\[
\operatorname{BIEXT}(P,Q;G)\xrightarrow{i}\operatorname{TORSBIRIG}(P,Q;G),\tag{2.9.1}
\]
其中第二项表示 $P\times Q$ 上（结构\Term{group}{群}为 $G_{P\times Q}$）的\Term{birigidification}{双刚化}\Term{torsor}{挠子}\Term{category}{范畴}。

利用 $P$ 和 $Q$ 的\Term{unit-section}{单位截面}，对于 $n$ 个分别与 $P$ 或 $Q$ \Term{isomorphic}{同构}的\Term{object}{对象}的任意乘积，都可得到从第 $i$ 个因子到该乘积的一个\Term{canonical}{典范}\Term{morphism}{态射} $\operatorname{inj}_i$。若给定从这个乘积到 $G$ 的一个\Term{morphism}{态射}，这样便可谈论该\Term{morphism}{态射}在第 $i$ 个因子上的\Term{restriction}{限制}。在此基础上，我们考虑以下乘积：
\[
\left\{\begin{gathered}
P\times Q,\\
P\times Q\times Q,\quad P\times Q\times Q\times Q,\\
P\times P\times Q,\quad P\times P\times P\times Q,\\
P\times P\times Q\times Q。
\end{gathered}\right.
\tag{2.9.2}
\]
\begin{SGAresult}{命题}{2.9.3}
假设从 (2.9.2) 中任一乘积到 $G$ 的每个\Term{morphism}{态射}，若在该乘积的每个因子上都\Term{trivial-morphism}{平凡}（各因子借助\Term{unit-section}{单位截面}\Term{embedding}{嵌入}乘积，如上所述），则它本身\Term{trivial-morphism}{平凡}。那么：

 a) \Term{canonical}{典范}\Term{functor}{函子} (2.9.1) 是\Term{fully-faithful}{充分忠实}的，而且 (2.9.1) 的两项都是\Term{rigid-category}{刚性范畴}。特别地，若 $E$ 是一个相对于 $(e_P,e_Q)$ \Term{birigidification}{双刚化}的 $G_{P\times Q}$-\Term{torsor}{挠子}，则 $E$ 上至多存在一个与其\Term{birigidification}{双刚化}\Term{torsor}{挠子}结构相容的\Term{biextension}{双扩张}结构。
\SourcePage{182}

 b) $E$ 上存在这样的\Term{biextension}{双扩张}结构，当且仅当 $E$ 上存在一个与其 $G_P$-\Term{torsor}{挠子}结构相容的、$Q_P$ 以 $G_P$ 为核的\Term{group-extension}{扩张}结构，并且 $E$ 上存在一个与其 $G_Q$-\Term{torsor}{挠子}结构相容的、$P_Q$ 以 $G_Q$ 为核的\Term{group-extension}{扩张}结构（参见 1.3.5 b)）。在这种情形，$E$ 上存在唯一的 $Q_P$ 以 $G_P$ 为核的\Term{group-extension}{扩张}结构，与 $G_P$-\Term{torsor}{挠子}结构及 $P\times e_Q$ 上给定的\Term{rigidification}{刚化}相容；也存在唯一的 $P_Q$ 以 $G_Q$ 为核的\Term{group-extension}{扩张}结构，与 $G_Q$-\Term{torsor}{挠子}结构及 $e_P\times Q$ 上给定的\Term{rigidification}{刚化}相容；而且这两个\Term{group-extension}{扩张}结构相容（2.1）。
\end{SGAresult}
\begin{proof}
a) 若 $E$ 是 $P\times Q$ 上一个\Term{birigidification}{双刚化}\Term{torsor}{挠子}，则其\Term{automorphism-group}{自同构群}等同于在两个部分因子上均\Term{trivial-morphism}{平凡}的\Term{morphism}{态射} $P\times Q\to G$ 所成的\Term{group}{群}，故由应用于乘积 $P\times Q$ 的假设，它只有恒等\Term{automorphism}{自同构}。这证明了\Term{rigidity}{刚性}断言。为证明若 $E,E'$ 是 $(P,Q)$ 以 $G$ 为核的\Term{biextension}{双扩张}，则\Term{birigidification}{双刚化}\Term{bitorsor}{双挠子}的任意\Term{morphism}{态射} $f:E\to E'$ 都是\Term{biextension}{双扩张}\Term{morphism}{态射}，只须分别将 1.3.5 a) 应用于 $E,E'$ 在 $P$ 和 $Q$ 上定义的两个\Term{group-extension}{扩张}结构（这将用到关于 (2.9.2) 中两个乘积 $P\times Q\times Q$ 和 $P\times P\times Q$ 的假设）。

b) 对\Term{birigidification}{双刚化}\Term{torsor}{挠子} $E$ 所述的条件，是它来自一个\Term{biextension}{双扩张}的必要条件，这很清楚；现证其充分性。首先，由 1.3.5 b) 得知，所述条件蕴含存在唯一的 $Q_P$ 以 $G_P$ 为核及 $P_Q$ 以 $G_P$ 为核的\Term{group-extension}{扩张}结构，它们与\Term{torsor}{挠子}结构及各\Term{rigidification}{刚化}相容。（这里用到关于 (2.9.2) 中乘积 $P\times Q\times Q\times Q$ 和 $P\times P\times P\times Q$ 的假设。）
\SourcePage{183}
尚须证明这两个\Term{group-extension}{扩张}结构相容，即图 (2.1.1) 交换。然而，该图中的两个复合都是 $P\times P\times Q\times Q$ 上的\Term{torsor}{挠子}\Term{isomorphism}{同构}，故它们相差一个\Term{morphism}{态射}
\[
u:P\times P\times Q\times Q\longrightarrow G。
\]
利用 $E$ 的各部分\Term{trivialization}{平凡化}分别为\Term{composition-law}{合成律} $\varphi$ 和 $\psi$ 的\Term{unit-section}{单位截面}这一事实，容易得出 $u$ 在这四个因子中的每一个上均\Term{trivial-morphism}{平凡}；于是，由应用于 (2.9.2) 中乘积 $P\times P\times Q\times Q$ 的假设，$u$ 本身\Term{trivial-morphism}{平凡}，证明完毕。
\end{proof}

\begin{SGAresult}{推论}{2.9.4}
设 $P,Q,G$ 是满足关系
\[
\underline{\operatorname{Hom}}_{\mathrm{pt}}(P,G)=e,\qquad
\underline{\operatorname{Hom}}_{\mathrm{pt}}(Q,G)=e\tag{2.9.4.1}
\]
的\Term{commutative-group}{交换群}（参见 1.3.7）。则 2.9.2 的前提条件成立，因此该命题的结论 a)、b) \Term{effective}{有效}。此外，若 $E$ 是相对于 $(e_P,e_Q)$ \Term{birigidification}{双刚化}的一个 $G_{P\times Q}$-\Term{torsor}{挠子}，则它具有一个与其给定结构相容的、$(P,Q)$ 以 $G$ 为核的\Term{biextension}{双扩张}结构，当且仅当它所定义的两个带点\Term{set}{集合}\Term{sheaf}{层}\Term{morphism}{态射}
\[
Q\longrightarrow R^1p_*(G_P),\qquad P\longrightarrow R^1q_*(G_Q)\tag{2.9.4.2}
\]
是\Term{group}{群}\Term{morphism}{态射}（$p,q$ 是结构\Term{morphism}{态射} $P\to e$、$Q\to e$）。
\end{SGAresult}
\begin{proof}
2.9.4 的第一个断言已在 1.3.7 中见到。对于第二个断言，所考虑的条件显然必要；现通过应用 2.9.2 b) 和 1.3.5 b) 证明其充分性。
\SourcePage{184}
注意，(2.9.4.2) 中第一个\Term{morphism}{态射}的\Term{multiplication}{乘法}性，意味着\Term{rigidification}{刚化}的 $G_P$-\Term{torsor}{挠子} $E$ 在 $P$ 上\Term{local-scheme}{局部}满足 1.3.5 b) 的条件。我们在 1.3.7 中已指出，这蕴含该条件在 $P$ 上整体成立，因此 $E$ 上存在一个与其\Term{rigidification}{刚化}的 $G_P$-\Term{torsor}{挠子}结构相容的、$Q_P$ 以 $G_P$ 为核的\Term{group-extension}{扩张}结构。同样可见，$E$ 上存在一个与其\Term{rigidification}{刚化}的 $G_Q$-\Term{torsor}{挠子}结构相容的、$P_Q$ 以 $G_Q$ 为核的\Term{group-extension}{扩张}结构。因此，结论由 2.9.2 b) 得出。
\end{proof}

\textbf{例 2.9.5。} 设 $S$ 为一个\Term{scheme}{概形}，$P,Q$ 是 $S$ 上两个\Term{abelian-scheme}{阿贝尔概形}，$G=\mathbf G_{m,S}$ 是 $S$ 上的\Term{multiplicative-group-scheme}{乘法群概形}。在 1.3.8 中已指出，这时条件 (2.9.4.1) 成立。另一方面，给定 $P\times Q$ 上一个\Term{birigidification}{双刚化}\Term{torsor}{挠子} $E$，\Term{equivalence}{等价}于给定 $P\times Q$ 上一个\Term{birigidification}{双刚化}\Term{invertible-module}{可逆模}。回忆，所谓 $P,Q$ 上的一个\Term{divisorial-correspondence}{除子对应}，是指这样一个\Term{birigidification}{双刚化}\Term{invertible-module}{可逆模}的\Term{isomorphism}{同构}类（例如参见 [1]）。在 fppf \Term{topology}{拓扑}中工作时，这里的\Term{morphism}{态射} (2.9.4.2) 是由\Term{divisorial-correspondence}{除子对应}定义的\Term{morphism}{态射}
\[
Q\longrightarrow\underline{\operatorname{Pic}}_{P/S},\qquad
P\longrightarrow\underline{\operatorname{Pic}}_{Q/S},\tag{2.9.5.1}
\]
由“\Term{theorem-of-the-square}{平方定理}”（同上）熟知，这些\Term{morphism}{态射}必定与\Term{group}{群}结构相容。所以 2.9.4 告诉我们，$(P,Q)$ 以 $G=\mathbf G_{m,S}$ 为核的\Term{biextension}{双扩张}\Term{category}{范畴}，与 $P\times Q$ 上\Term{birigidification}{双刚化}\Term{invertible-module}{可逆模}的\Term{category}{范畴}\Term{equivalence-of-categories}{等价}；因此它是一个\Term{rigid-category}{刚性范畴}，其\Term{isomorphism}{同构}类\Term{set}{集合}等同于 $(P,Q)$ 上的\Term{divisorial-correspondence}{除子对应}\Term{set}{集合}。
\SourcePage{185}
此外，这个等同显然保持两个\Term{set}{集合}上的自然\Term{group}{群}律（第一个\Term{group}{群}律定义于 (2.5.2.2)）。

\textbf{注记 2.9.6。} 容易验证，(2.9.4.2) 中第一个\Term{morphism}{态射}的\Term{additivity}{加性}，\Term{equivalence}{等价}于说 (2.9.4.2) 中第二个\Term{morphism}{态射}经由子\Term{sheaf}{层} $\underline{\operatorname{Ext}}^1(Q,G)$ \Term{resolution}{分解}（注意，由 1.3.5，\Term{canonical}{典范}\Term{morphism}{态射}
\[
\underline{\operatorname{Ext}}^1(Q,G)\longrightarrow R^1q_*(G_Q)
\]
是\Term{injective-morphism}{单射}）。同样，(2.9.4.2) 中第二个\Term{morphism}{态射}的\Term{additivity}{加性}，意味着第一个\Term{morphism}{态射}经由 $\underline{\operatorname{Ext}}^1(P,G)$ \Term{resolution}{分解}。此外，由此得到\Term{group}{群}\Term{isomorphism}{同构}
\[
\operatorname{Biext}(P,Q;G)\simeq\operatorname{Hom}\bigl(Q,\underline{\operatorname{Ext}}^1(P,G)\bigr)
\simeq\operatorname{Hom}\bigl(P,\underline{\operatorname{Ext}}^1(Q,G)\bigr)。\tag{2.9.6.1}
\]
在 2.9.5 考虑的特殊情形，这些\Term{isomorphism}{同构}化为熟知的\Term{isomorphism}{同构}
\[
\operatorname{Corr}(P,Q)\simeq\operatorname{Hom}(P,Q')\simeq\operatorname{Hom}(Q,P'),\tag{2.9.6.2}
\]
其中 $P'=\underline{\operatorname{Ext}}^1(P,\mathbf G_m)$ 和 $Q'=\underline{\operatorname{Ext}}^1(Q,\mathbf G_m)$ 是 $P$ 和 $Q$ 的\Term{dual-abelian-scheme}{对偶阿贝尔概形}，分别等同于 $\operatorname{Pic}^\circ_{P/S}$ 和 $\operatorname{Pic}^\circ_{Q/S}$。当然，在前面两个公式中，$\operatorname{Hom}$ 表示\Term{group}{群}\Term{homomorphism}{同态}所成的\Term{group}{群}。

\textbf{2.10。\Term{biextension}{双扩张}概念的若干推广。} 为求周全，我们将它们列出，不过以后并不需要使用。

\textbf{2.10.1。} 不必假设 $P,Q$ 交换，只假设 $G$ 交换，也可以定义 $(P,Q)$ 以 $G$ 为核的\Term{biextension}{双扩张}概念：它是一个 $G_{P\times Q}$-\Term{torsor}{挠子}，
\SourcePage{186}
赋有 $Q_P$ 以 $G_P$ 为核和 $P_Q$ 以 $G_Q$ 为核的\Term{central-extension}{中心扩张}结构，这两个结构与\Term{torsor}{挠子}结构相容，并且彼此相容，即使图 (2.1.1) 交换。前面的大部分讨论经适当修改后，对这个更\Term{generic}{一般}的概念仍然\Term{effective}{有效}。须注意，即使 $P,Q$ 也交换，上述概念仍比 2.1 的概念更\Term{generic}{一般}；2.1 对应于 $E$ 的两个底层\Term{group-extension}{扩张}结构都是相对\Term{commutative-group}{交换群}结构的情形。还须注意，在非交换情形，$(P,Q)$ 以 $G$ 为核的\Term{biextension}{双扩张}\Term{category}{范畴}仍\Term{additivity}{加性地}依赖于 $G$，却不再\Term{additivity}{加性地}依赖于 $P,Q$。

\textbf{2.10.2。} 照搬 2.1，便可（对给定整数 $n\geq1$）定义一列 $n$ 个\Term{commutative-group}{交换群} $(P_1,P_2,\ldots,P_n)$ 以一个\Term{commutative-group}{交换群} $G$ 为核的\Term{n-extension}{$n$ 重扩张}概念。它对应于在一个 $G_{P_1\times\cdots\times P_n}$-\Term{torsor}{挠子} $E$ 上，给定 $n$ 个以 $n-1$ 个因子的各部分乘积为底的部分\Term{group-extension}{扩张}结构，并对 $1$ 至 $n$ 间每对不同指标，有一个形如 (2.1.1) 的相容性关系。

当 $n\geq3$，且 $P_i$ 是一个\Term{scheme}{概形} $S$ 上的\Term{abelian-scheme}{阿贝尔概形}，$G$ 是 $\mathbf G_m$ 上的\Term{multiplicative-group}{乘法群}时（比较 2.9.2），可发现 $(P_1,\ldots,P_n)$ 以 $G$ 为核的任意\Term{n-extension}{$n$ 重扩张}都是\Term{trivial-extension}{平凡}的。

\textbf{2.10.3。} 2.1 的\Term{biextension}{双扩张}概念，可以看作\Term{commutative-group}{交换群}的\Term{commutative-group-extension}{交换扩张}概念，即 $\mathbf Z$-\Term{module}{模}的\Term{group-extension}{扩张}概念的一个变体。给定 $\mathcal T$ 的一个\Term{ring}{环} $A$（为简单起见假设交换），以及三个 $A$-\Term{module}{模} $P,Q,G$，可以
\SourcePage{187}
更\Term{generic}{一般}地定义 $(P,Q)$ 以 $G$ 为核的 $A$-\Term{module}{模}\Term{biextension}{双扩张}概念，它是 $A$-\Term{module}{模}\Term{module-extension}{扩张}概念的一个变体。这样的\Term{biextension}{双扩张}定义为 2.1 意义下的一个\Term{biextension}{双扩张}，但把底层的各部分\Term{group-extension}{扩张}结构充实为 $A_P$-\Term{module}{模}及 $A_Q$-\Term{module}{模}的\Term{module-extension}{扩张}结构（与 $Q_P,G_P$ 和 $P_Q,G_Q$ 上已有的相对\Term{module}{模}结构相容）。$E$ 上的两个相对\Term{module}{模}结构，\Term{solvable}{可解}释为 $A$ 的底层\Term{multiplication}{乘法}\Term{monoid}{幺半群}在 $E$ 的底层\Term{set}{集合}\Term{sheaf}{层}上的两个\Term{group-action}{作用}；我们假设这两个\Term{group-action}{作用}彼此交换。因此，这个相容性是图 (2.1.1) 的\Term{commutativity}{交换性}所表达的相容性之外的附加条件，它由一个类似的图的\Term{commutativity}{交换性}表达，留给读者写出。最后，假设 $A$ 在 $E$ 上的第一个\Term{group-action}{作用}，不仅对 $E$ 的第一个\Term{composition-law}{合成律}，也对 $E$ 的第二个\Term{composition-law}{合成律}满足\Term{distributivity}{分配律}；对于 $A$ 在 $E$ 上的第二个\Term{group-action}{作用}也作同样假设。

还应看到，1.1.6 中所述的对应也有一个 $A$-\Term{linear}{线性}变体：给定 $A$-\Term{module}{模} $P$ 以另一个 $A$-\Term{module}{模} $G$ 为核的一个\Term{module-extension}{扩张}，\Term{equivalence}{等价}于给定一族以“点” $p\in P(S)$ 为参数的 $G_p$-\Term{torsor}{挠子} $E_p$（$S\in\operatorname{ob}\mathcal T$ 为可变\Term{object}{对象}），并要求 $E_p$ 随 $p$ 以 $A$-\Term{linear}{线性}方式变化（而不仅像 1.1.6 中那样\Term{additivity}{加性地}变化）。具体地说，除\Term{isomorphism}{同构} $\varphi_{p,p'}$（1.1.3.1）外，还须给定\Term{isomorphism}{同构}
\[
\mu_{a,p}:{}^a(P_p)\xrightarrow{\sim}P_{ap}\qquad(p\in P(S),\ a\in A(S)),
\]
其中第一项表示 $G_S$-\Term{torsor}{挠子} $P_p$ 经由\Term{extension-of-scalars}{标量扩张}
\SourcePage{188}
$g\mapsto ag:G_S\to G_S$ 所得的像。至于须对这些数据 $\varphi,\mu$ 施加什么条件，才能使之确实对应于 $P$ 以 $G$ 为核的一个\Term{module}{模}\Term{module-extension}{扩张} $E$，留给读者具体写出（参见 3.5.4 b)）。
